\section{Conclusion}
\label{sec:conclusion}

We integrated a feed-forward Gaussian predictor into a dense SLAM system and
asked what has to change for the \emph{fused map}, rather than the individual
prediction, to improve. Adaptation must be confined to the Gaussian head:
perturbing a backbone whose downstream head was trained against frozen
features collapses reconstruction by $49\%$, whereas head-only training helps
on five families. A training-free lever --- fading opacity at injection in
proportion to pointmap confidence --- helps in $10$ of $12$ online cells with no
harmful cell, and cannot be replaced by a training-time penalty because
crowding is a map-level property that pair training never observes. Both
results are explained by one mechanism: opacity in this architecture is a
brightness-trim dial against the black-background rasteriser, confirmed by a
pre-registered opposite-sign prediction that held on $5/5$ families.

Evaluated through a single protocol against four locally-built systems, our
method leads Photo-SLAM on all eight Replica scenes in PSNR ($+3.44$\dB{}
mean) while winning LPIPS on $5/8$, ties MonoGS on TUM at matched compute
while losing LPIPS there, and inherits MASt3R-SLAM-class tracking robustness
that the two GS-SLAM baselines do not match. It does so with a map
$25$--$100\times$ larger, and we show that this size is constitutive rather
than redundant.

Two methodological points seem to us more transferable than the system itself.
First, GS-SLAM rendering numbers in the literature and ours differ by
approximately $8.7$\dB{} of \emph{protocol} --- measured twice, on two systems,
by two independent methods --- so cross-paper rendering comparisons in this
area should be treated as unreliable unless the protocols are stated and
matched. Second, in the course of this work three separate results that
initially favoured us changed when checked one level deeper: an inverted pose
convention, a misread self-consistency test, and a small-scene subset that
overstated our margin by $1.85$\dB{}. We adopted the rule that a result which
flatters the authors earns one more verification step than one that does not,
and we recommend it.
